\chapter{Как машина учится двигаться}
\chaptermark{Как машина учится двигаться}
\label{sec:motorlearning}

\section{Три учителя}

В главе~\ref{sec:muscles} мы видели, как машина двигает мышцами, и остановились на гипотезе: быстрые и точные движения требуют внутренней модели тела, которая предсказывает результат команды~\cite{WolpertGhahramani2000}. Откуда берётся эта модель? Её надо выучить, и выучить не так, как мы учились до сих пор.

В книге уже было два учителя. Первый --- \emph{никакой}: правило Хебба (глава~\ref{sec:dofa}) усиливает то, что часто бывает вместе, и сжимает входы. Второй --- \emph{награда}: \nt{DOPA} рассылает всем одно число, «лучше или хуже, чем ожидалось». Движению нужен третий учитель --- \emph{ошибка}. Рука промахнулась мимо чашки на два сантиметра влево: это не «хуже», это вектор, который говорит каждому выходу, куда и насколько он ошибся. Инженер назвал бы это обучением с учителем.

Разница между вторым и третьим учителем --- это разница между одним проводом на всех и отдельным проводом на каждый выход. В утверждении~\ref{prop:broadcastcost} обучение по одному общему числу замедлялось с каждым нейроном, чей шум влиял на результат. Если же каждый выход $i$ получает свою ошибку $e_i$, веса можно менять по простейшему правилу:
\begin{empheq}[box=\keybox]{equation}
\frac{\dd w_{ij}}{\dd t} \;=\; -\eta\,e_i(t)\,x_j(t) .
\label{eq:lms}
\end{empheq}
Это правило наименьших квадратов, давно известное в технике адаптивных фильтров~\cite{WidrowHoff1960}: вес меняется пропорционально произведению своего входа на ошибку своего выхода и в среднем идёт по градиенту квадрата ошибки. Никакого шума-поиска, как в главе~\ref{sec:dofa}, не нужно: направление поправки приходит прямо по проводу ошибки. И скорость не падает с числом выходов, потому что чужие ошибки к синапсу не приходят.

\begin{keyblock}
\begin{proposition}[Три учителя]
\label{prop:teachers}
Правило Хебба учит без учителя тому, что часто; сигнал \nt{DOPA} учит одним числом на всю сеть тому, что выгодно, и медленнее с каждым нейроном; провод ошибки к каждому выходу учит по правилу~\eqref{eq:lms} тому, что точно, и скорость не зависит от числа выходов. Цена третьего способа --- отдельный провод на каждый выход и кто-то, кто знает правильный ответ.
\end{proposition}
\end{keyblock}

\section{Провод ошибки}

Кто может знать правильный ответ для движения? Сами датчики. Если рука ушла влево от цели, об этом сообщают глаз и датчики растяжения из главы~\ref{sec:sensors}. Нужна схема, которая доставит эту ошибку к нужным выходам.

В мозге есть область, которая, по-видимому, устроена ровно под это~\cite{Marr1969,Albus1971}. Её схема необычно правильная, почти кристаллическая, и в ней легко узнать правило~\eqref{eq:lms} (рис.~\ref{fig:errorwire}).

\emph{Расширитель входа.} Сигналы о команде и о состоянии тела приходят на огромный слой маленьких \nt{GLUT}-нейронов. Их около семидесяти миллиардов --- больше, чем во всём остальном мозге, вместе взятом~\cite{Azevedo2009}. Каждый из них смешивает несколько входов, и сигнал раскладывается в очень длинный вектор. Инженеру этот приём знаком: если поднять размерность входа, в новом пространстве почти любую нужную зависимость можно получить простой взвешенной суммой~\cite{Cover1965}.

\emph{Выходные нейроны.} Взвешенную сумму считают около тридцати миллионов больших выходных нейронов~\cite{Andersen1992cereb}, у каждого --- порядка ста тысяч входов от расширителя. И это \nt{GABA}-нейроны: результат обучения передаётся дальше не возбуждением, а торможением, как в запрете главы~\ref{sec:gaba}.

\emph{Провод ошибки.} Каждый выходной нейрон получает ещё один вход, особый: ровно один \nt{GLUT}-провод, который срабатывает редко, около раза в секунду, но каждый раз вызывает в выходном нейроне мощную вспышку~\cite{Ito2001}. Это и есть $e_i$: вероятность вспышки зависит от направления ошибки движения~\cite{Herzfeld2018}. Совпадение вспышки ошибки с активностью входа ослабляет этот вход --- в точности член $-\eta\,e_i x_j$ в~\eqref{eq:lms}.

\begin{figure}[t]
\centering
\begin{tikzpicture}[>=Stealth,
  gran/.style={circle,draw,thick,fill=blue!6,minimum size=0.32cm,inner sep=0pt},
  outn/.style={circle,draw,very thick,fill=red!10,minimum size=1.1cm,inner sep=0pt,font=\scriptsize},
  inh/.style={-{Bar[width=7pt]},very thick,red!70!black}]
% inputs
\foreach \y/\lab in {1.6/{команда},0.4/{состояние тела}}{
  \draw[->,thick,blue!60!black] (-0.6,\y) node[left,font=\scriptsize,black] {\lab} -- (0.35,\y);}
% expansion layer
\foreach \i in {0,...,11}{\node[gran] (g\i) at ({0.8+0.28*mod(\i,2)},{-0.3+0.23*\i}) {};}
\foreach \i in {0,...,11}{\pgfmathsetmacro{\yy}{mod(\i,3)>0 ? 1.6 : 0.4}\draw[gray!60!black] (0.35,\yy) -- (g\i);}
\node[font=\scriptsize,align=center] at (0.95,-1.05) {расширитель\\$\sim7\cdot10^{10}$ \nt{GLUT}};
% parallel wires to the output neuron
\foreach \i in {0,...,11}{\draw[->,gray!70!black] (g\i.east) -- (4.1,{1.0+0.07*(\i-5.5)});}
\node[outn] (P) at (4.6,1.0) {\nt{GABA}};
\node[font=\scriptsize,align=center,above=2pt] at (P.north) {выходной нейрон\\$\sim10^5$ входов $x_j$, веса $w_{ij}$};
\draw[inh] (P) -- (6.8,1.0) node[right,font=\scriptsize,align=left] {поправка\\к движению};
% error wire
\draw[->,very thick,red!70!black,densely dashed] (4.6,-1.4) node[below,font=\scriptsize,black,align=center] {один провод ошибки $e_i$\\\nt{GLUT}, $\sim1$ раз/с} -- (P.south);
\end{tikzpicture}
\caption{Схема обучения с учителем, как её, по-видимому, реализует мозг. Команда и состояние тела раскладываются огромным слоем \nt{GLUT}-нейронов в длинный вектор $x_j$; выходной \nt{GABA}-нейрон складывает его с весами $w_{ij}$. Единственный провод ошибки приносит $e_i$; совпадение ошибки с активным входом ослабляет вес этого входа, как в правиле~\eqref{eq:lms}.}
\label{fig:errorwire}
\end{figure}

Одна трудность знакома нам по главе~\ref{sec:dofa}: ошибка приходит позже, чем вход, который к ней привёл. Промах виден через десятки-сотни миллисекунд после команды. Значит, синапс должен помнить, что был активен, --- нужен след, как в правиле~\eqref{eq:threefactor}. И длина этого следа, судя по опытам, подстроена под задержку обратной связи в каждой задаче: там, где ошибка приходит через сто миллисекунд, сильнее всего ослабляется вход, бывший активным за сто миллисекунд до неё~\cite{Suvrathan2016}.

\begin{keyblock}
\begin{proposition}[Провод ошибки]
\label{prop:errorwire}
В схеме, изображённой на рис.~\ref{fig:errorwire}, каждый выходной нейрон получает ровно один провод ошибки, и совпадение ошибки с входом ослабляет этот вход. Это правило наименьших квадратов~\eqref{eq:lms}, применённое к входу, который расширен до огромной размерности. Устройство схемы и ослабление при совпадении измерены; что вся схема работает как обучение с учителем, --- ведущая гипотеза.
\end{proposition}
\end{keyblock}

\section{Внутренняя модель как адаптивный фильтр}

Чему учится такая схема? Самый естественный ответ --- предсказанию. Пусть на вход приходит команда $u(t)$, пропущенная через набор фильтров с разными постоянными времени, как в расширителе: $x_j(t)=(g_j*u)(t)$. Выход --- их взвешенная сумма, предсказание того, что сообщат датчики:
\begin{empheq}[box=\keybox]{equation}
\hat y(t) \;=\; \sum_j w_j\,(g_j*u)(t), \qquad e(t) \;=\; \hat y(t) - y(t) ,
\label{eq:adaptivefilter}
\end{empheq}
а веса учатся по~\eqref{eq:lms}. Это \emph{адаптивный фильтр}: схема, которая сама подбирает свою передаточную функцию так, чтобы повторить неизвестную систему~\cite{Fujita1982}. Здесь неизвестная система --- собственное тело с его массой, упругостью и задержками. Выученный фильтр и есть внутренняя модель главы~\ref{sec:muscles}: он говорит, что сообщат датчики, не дожидаясь их.

Такая модель полезна дважды. Во-первых, её предсказание можно подать вместо запаздывающей обратной связи, и тогда условие устойчивости~\eqref{eq:delaystable} перестаёт связывать руки: исправлять можно быстро, по модели. Во-вторых, разница между предсказанным и полученным --- это сигнал о \emph{неожиданном}. Всё, что машина сделала сама, модель предсказала и вычла; остаётся то, что сделал мир. Поэтому, по одной из гипотез, мы и не можем пощекотать сами себя~\cite{Blakemore1998}.

\section{Когда мир меняется: быстрое и медленное}

Проверим схему опытом, который легко поставить. Человек тянется к цели, а мир неожиданно меняется: например, на руку действует боковая сила. Первые движения промахиваются, потом промах уменьшается. Если силу убрать, рука сначала промахивается в \emph{другую} сторону --- модель выучила силу и продолжает её компенсировать. Это прямое свидетельство того, что выучена модель, а не просто «стало получаться».

Опыты показывают и более тонкую вещь: учатся одновременно два процесса --- быстрый, который быстро учится и быстро забывает, и медленный, который учится медленно и помнит долго~\cite{Smith2006}. Поправки обновляются после каждого движения, по событию:
\begin{empheq}[box=\keybox]{equation}
e = p - (x_f + x_s), \qquad x_f \leftarrow A_f\,x_f + B_f\,e, \qquad x_s \leftarrow A_s\,x_s + B_s\,e ,
\label{eq:tworate}
\end{empheq}
где $p$ --- возмущение, $x_f$ и $x_s$ --- выученные поправки двух процессов, $A$ --- насколько поправка сохраняется до следующего движения, $B$ --- насколько сильно она учится на ошибке. У быстрого процесса $B$ велико, а $A$ заметно меньше единицы; у медленного --- наоборот.

\begin{figure}[t]
\centering
\begin{tikzpicture}[>=Stealth,x=0.034cm,y=1.9cm]
\draw[->,gray!70!black] (0,0) -- (355,0) node[right,font=\small,black] {движение};
\draw[->,gray!70!black] (0,-1.15) -- (0,1.25);
\foreach \v in {-1,1}{\draw[gray!60!black] (-3,\v) -- (3,\v); \node[left,font=\scriptsize] at (0,\v) {$\v$};}
\foreach \n in {100,200,300}{\draw[gray!60!black] (\n,-0.04) -- (\n,0.04); \node[below,font=\scriptsize] at (\n,0) {\n};}
\draw[thin,gray!70!black] plot file {figures/tworate_p.dat};
\draw[thick,orange!85!black] plot file {figures/tworate_fast.dat};
\draw[thick,blue!60!black] plot file {figures/tworate_slow.dat};
\draw[very thick,black] plot file {figures/tworate_net.dat};
\node[font=\scriptsize,gray!50!black] at (120,1.12) {возмущение $p$};
\node[font=\scriptsize,black,anchor=west] at (238,0.52) {сумма: возврат};
\node[font=\scriptsize,blue!60!black,anchor=west] at (150,0.39) {медленный $x_s$};
\node[font=\scriptsize,orange!85!black,anchor=west] at (60,0.08) {быстрый $x_f$};
\draw[densely dashed,gray!60!black] (226,-1.15) -- (226,1.25);
\node[font=\scriptsize,align=center,anchor=south west] at (228,-1.1) {ошибок\\больше нет};
\end{tikzpicture}
\caption{Два процесса обучения по модели~\eqref{eq:tworate}, $A_f=0.92$, $B_f=0.1$, $A_s=0.996$, $B_s=0.02$. Двести движений с возмущением $p=1$, затем шесть с обратным, $p=-1$, пока сумма не вернётся к нулю. После штриховой линии ошибки перестают поступать, и сумма сама возвращается к старой поправке: быстрый процесс забыл обратное возмущение, медленный помнит прямое.}
\label{fig:tworate}
\end{figure}

Модель~\eqref{eq:tworate} объясняет опыт, который трудно понять без неё (рис.~\ref{fig:tworate}). В нашем расчёте за двести движений с возмущением сумма поправок доходит до $0.84$, причём большую часть держит медленный процесс, $0.63$. Затем шесть движений с обратным возмущением возвращают сумму к нулю: быстрый процесс ушёл в минус, $-0.53$, медленный ещё в плюсе, $0.46$. Если теперь перестать сообщать ошибки, быстрый процесс забывает свою поправку за десятки движений, медленный --- гораздо дольше, и сумма \emph{сама} поднимается до $0.37$ за сорок движений: старый навык возвращается без всякой тренировки. Такое спонтанное восстановление измерено у людей; модель с одним процессом его дать не может --- без ошибок она просто затухает к нулю.

Зачем два процесса? Правдоподобный инженерный ответ даёт оптимальный фильтр главы~\ref{sec:acet}. Тело меняется по разным причинам с разной скоростью: мышцы устают за минуты, растут и слабеют за месяцы. Если помехи бывают быстрыми и медленными, наилучшая оценка состоит из двух фильтров с разными постоянными времени, и каждый ловит свою~\cite{Kording2007}. Быстрый процесс --- это временная поправка, «рука устала»; медленный --- «рука стала другой». Это пока гипотеза, но она объясняет, почему навык, выученный за один день, частично теряется к следующему и почему второй раз он учится быстрее.

\section{Итог}

\begin{table}[t]
\centering
\footnotesize
\setlength{\tabcolsep}{4pt}
\renewcommand{\arraystretch}{1.15}
\begin{tabular}{>{\raggedright}p{3.0cm}>{\raggedright}p{4.6cm}>{\raggedright\arraybackslash}p{5.2cm}}
\toprule
Что делает схема & Как & Что известно \\
\midrule
учится с учителем & провод ошибки к каждому выходу, правило~\eqref{eq:lms} & устройство схемы и ослабление при совпадении измерены; обучение с учителем --- ведущая гипотеза \\
расширяет вход & семьдесят миллиардов \nt{GLUT}-нейронов & число нейронов измерено; роль расширения --- теория \\
учитывает задержку ошибки & след, подстроенный под задержку & измерено \\
строит внутреннюю модель & адаптивный фильтр~\eqref{eq:adaptivefilter} & хорошо обоснованная гипотеза \\
учится двумя темпами & быстрый и медленный процессы~\eqref{eq:tworate} & последействие и спонтанный возврат измерены; два процесса --- модель \\
\bottomrule
\end{tabular}
\caption{Как машина учится двигаться.}
\label{tab:motorlearning}
\end{table}

Теперь у машины три учителя (таблица~\ref{tab:motorlearning}). Хебб сжимает то, что часто; \nt{DOPA} одним числом учит выбирать выгодное; провод ошибки учит точности, и учит быстро, потому что к каждому выходу идёт своя ошибка. Мозг, по-видимому, пользуется всеми тремя, и каждым там, где он дешевле всего: широковещательным сигналом --- для немногих решений, отдельными проводами --- для точной подстройки тела, где правильный ответ сообщают сами датчики.

\section{Задачи}

\begin{exercise}[\lvlA]
\label{ex:15:1}
\emph{Ёмкость схемы с проводом ошибки.} В схеме $3\cdot10^7$ выходных нейронов по $10^5$ входов, у каждого свой провод ошибки ($1$~имп/с). Нейрон с $n$ входами выучивает любую разметку $\approx2n$ векторов~\cite{Cover1965}. (а)~Сколько ассоциаций хранит схема? (б)~Оцените по~\eqref{eq:spikeentropy} при $\Delta t=10\ms$ наименьшее время заполнения ёмкости нейрона.
\end{exercise}

\begin{exercise}[\lvlB]
\label{ex:15:2}
\emph{Скорость правила наименьших квадратов.} Моды правила~\eqref{eq:lms} затухают со скоростями $\eta\lambda_k(C)$. Для пяти фильтров~\eqref{eq:adaptivefilter} на белом шуме $C_{jk}=2\sqrt{\tau_j\tau_k}/(\tau_j+\tau_k)$. Найдите отношение самой быстрой и самой медленной скоростей, если $\tau_j$ растут в $2$ раза ($10$--$160\ms$) и в $4$ раза.
\end{exercise}

\begin{exercise}[\lvlB]
\label{ex:15:3}
\emph{Анализ двух процессов.} Запишите модель~\eqref{eq:tworate} с параметрами рис.~\ref{fig:tworate} как $\mathbf x_{n+1}=M\mathbf x_n+\mathbf b\,p$. (а)~Найдите по собственным числам $M$ постоянные времени в движениях. (б)~Найдите установившиеся $x_f$, $x_s$ и их сумму при постоянном $p=1$.
\end{exercise}

\begin{exercise}[\lvlB]
\label{ex:15:4}
\emph{Моделирование: второй раз быстрее.} По модели~\eqref{eq:tworate} с параметрами из \texttt{models/\allowbreak motor\_learning.py} найдите число движений до $x_f+x_s\ge0.8$ при $p=1$: (а)~у необученной машины; (б)~после $200$ движений с $p=1$ и обратного возмущения до нуля суммы, как на рис.~\ref{fig:tworate}. Может ли так ускориться модель с одним процессом?
\end{exercise}

\begin{exercise}[\lvlB]
\label{ex:15:5}
\emph{Регулятор с внутренней моделью.} Объект $\dd x/\dd t=ku$ виден с задержкой $d$; регулятор $u=-G\hat x$ предсказывает $\hat x(t)=x(t-d)+\hat k\int_{t-d}^{t}u\,\dd s$. (а)~При $\hat k=k=1$, $d=150\ms$ найдите $G$ для постоянной времени $20\ms$ и сравните с пределом~\eqref{eq:delaystable}. (б)~Устойчива ли схема при $\hat k=0.4k$?
\end{exercise}

\begin{exercise}[\lvlB]
\label{ex:15:6}
\emph{Адаптивный фильтр учит мышцу.} Объект --- подёргивание~\eqref{eq:twitch} единичной площади, $T_c=50\ms$; фильтры~\eqref{eq:adaptivefilter} --- RC-цепи с $\tau_j=12.5$, $25$, $50$, $100$, $200\ms$; вход --- новое нормальное число каждые $10\ms$. За сколько секунд правило~\eqref{eq:lms} при $\eta=50$ и $200$ снижает ошибку до $0.5\%$? Почему веса и через $300$~с далеки от наилучших?
\end{exercise}

\begin{exercise}[\lvlC]
\label{ex:15:7}
\emph{Два процесса как оптимальный фильтр.} Возмущение --- сумма процессов $p^{f,s}_{n+1}=A_{f,s}\,p^{f,s}_n+w^{f,s}_n$ с дисперсиями шумов $q_f$, $q_s$, ошибка наблюдается с шумом дисперсии $R$; фильтр Калмана даёт~\eqref{eq:tworate}. При каких $q_f/R$, $q_s/R$ из $A_f=0.92$, $A_s=0.996$ выходит $B_f=0.1$, $B_s=0.02$? Как изменятся $B_f$, $B_s$, если учетверить $q_s$?
\end{exercise}
